% =====================================================================================
% Chapter 8 · Incrementality — media mix modelling
%
% Built from notebooks/06_incrementality_mmm.ipynb (kernel cmp-legacy, CMP_FAST=0).
% NOT ONE NUMBER IN THIS FILE IS TYPED BY HAND: every figure in the prose is a macro from
% build/macros.tex, generated by cmp.report from the executed notebook. The only literals
% below are (i) the constants of the data-generating model in §8.2, which are *definitions*
% rather than results; (ii) the decision-rule thresholds (0.9 / 0.8 / 0.5), the interval
% level (90 %) and the assumed 30 % gross margin, which are *conventions* stated in the
% text; (iii) the prior scales of §8.5, which define the model rather than measure anything;
% and (iv) the notebook's fast-mode series length (104 weeks), a configuration constant of
% the lab, cited once in §8.7 and never used as a result.
% =====================================================================================
\chapter{Incrementality: media mix modelling}
\label{ch:mmm}

\begin{quote}\small
\emph{A media mix model is handed \nbSixNWeeks{} weeks of spend and sales, the correct causal
graph, and the exact functional form that generated the data. It returns a confident answer, and
the answer is backwards. Television, which truly moved \eur{\nbSixTrueTv} of sales, is credited
\eur{\nbSixDefTv}; brand search, which truly moved \eur{\nbSixTrueBs}, is credited
\eur{\nbSixDefBs} --- so the model ranks brand search \emph{above} television, and it does so with
$\hat R = \nbSixDefRhat$, \nbSixDefDiv{} divergences, and a posterior predictive check that puts
\pc{\nbSixPpcCovDef} of weeks inside its \pc{90} band. Nothing in its own output warns anyone. The
repair is not a better prior; it is an experiment. One geo lift test (Chapter~9), translated
into a prior on brand search's saturation ceiling, restores the ranking ---
\eur{\nbSixCalTv} against \eur{\nbSixCalBs}. \Cref{sec:mmm:compare} then reports the
uncomfortable scoreline: on the decision-relevant quantity the \emph{classical} estimator beat the
default Bayesian one, and Bayes, out of the box, made things worse.}
\end{quote}

% =====================================================================================
\section{The decision}
\label{sec:mmm:decision}

A retailer spends on two channels. Television costs \eur{\nbSixSpendTv} over the planning window;
brand search --- the paid ads that appear when somebody types the company's own name into a search
engine --- costs \eur{\nbSixSpendBs}. The click-tracking dashboard is unambiguous: brand search
converts extraordinarily well, at a cost per acquisition television cannot approach. The proposal
on the table is to move budget out of television and into brand search.

The proposal is a catastrophe, and the reason is one sentence long. \textbf{People who are already
going to buy search for the brand.} The click that the dashboard credits to brand search was
preceded by an intention the channel did not create --- an intention created, perhaps, by the
television campaign the dashboard is about to defund, or simply by the season. Brand-search spend
is correlated with a sale it did not cause, and last-click attribution books that correlation as
credit. This is not a stylised worry. \citet{blake2015} switched eBay's paid search advertising
off across matched US markets and found the return on brand-keyword spend to be indistinguishable
from zero: the traffic simply arrived through the organic link instead. The channel's measured
return was, essentially, the cost of buying clicks the firm was already getting free.

\emph{Media mix modelling} is the standard defence. Instead of tracking clicks, regress weekly
sales on weekly spend per channel, plus controls, and read each channel's contribution off the
fitted model \citep{little1979, chan2017}. It is top-down where attribution is bottom-up; it sees
television, which sets no cookie; and --- decisively --- it is the only method most firms can
apply to the data they already have. It is, for those reasons, the most widely used method in this
book.

It is also the most dangerous, and for one reason that no amount of statistical sophistication
removes: \textbf{the spend was chosen, not randomized}. Nobody flipped a coin over the media plan.
Budgets were set by people who knew what season it was, what the competition was doing, and what
the sales forecast said. The regression therefore has no design behind it --- and it will produce
a confident, well-diagnosed, fully converged answer regardless. This chapter builds that model
properly, watches it fail, diagnoses the failure exactly, and repairs it with the one thing that
can: an experiment.

Two quantities must be separated before anything else, because the entire budget question turns on
the difference and the industry routinely conflates them.

\begin{itemize}[leftmargin=*]
  \item \textbf{Average ROAS} --- incremental revenue divided by the money already spent. It is a
        report card on the past.
  \item \textbf{Marginal ROAS} --- what the \emph{next} euro into the channel would buy, at
        today's spend. It is the only number a budget can be set from.
\end{itemize}

On a saturating channel the second is always below the first, and it can be far below: a channel
with a splendid average ROAS may be worth almost nothing at the margin because it is already
saturated. \Cref{sec:mmm:euro} measures exactly that gap in this data, and the two figures differ
by a factor of \nbSixAvgOverSat{}.

% =====================================================================================
\section{The data-generating model}
\label{sec:mmm:dgm}

A causal method cannot be graded on real data, because on real data the incremental contribution
of a channel is precisely what is missing. So the model is graded against a world whose answer is
known by construction --- and, so that the grading is as unforgiving as possible, that world is
built out of the media mix model's own functional form. Whatever goes wrong below cannot be blamed
on a mis-specified curve.

The simulator produces \nbSixNWeeks{} weeks of a weekly sales series. Write
$s_t = \sin(2\pi t / 52)$ for the seasonal index. Spend in each channel is drawn as

\begin{align}
  x^{\text{tv}}_t
    &= \max\big\{1,\ 40 + 8\,s_t + \varepsilon^{\text{tv}}_t\big\},
       &\varepsilon^{\text{tv}}_t &\sim \mathcal N(0,\, 8^2), \label{eq:mmm:xtv}\\
  x^{\text{bs}}_t
    &= \max\big\{1,\ 25 + 11\,s_t + \varepsilon^{\text{bs}}_t\big\},
       &\varepsilon^{\text{bs}}_t &\sim \mathcal N(0,\, 7^2). \label{eq:mmm:xbs}
\end{align}

\Cref{eq:mmm:xbs} is the whole trap, planted in one line. Brand-search spend is \emph{built} from
the seasonal index --- exactly as it is in the real world, where a firm bids harder on its own name
when demand is high and the search volume that triggers the spend \emph{is} the demand. Its
correlation with the season comes out at \nbSixCorrBsSeason{}, against television's
\nbSixCorrTvSeason{} (\cref{fig:nb06:confound}).

Spend becomes exposure through two transforms, which are the two ideas that make a media mix model
more than a regression.

\begin{description}[leftmargin=*, style=nextline]
  \item[Adstock, or carry-over.] An advertisement does not stop working the week it airs. Geometric
        adstock accumulates spend with a decay rate $\lambda \in [0, 1)$,
        %
        \begin{equation}
          a_t(x;\lambda) \;=\; x_t \;+\; \lambda\, a_{t-1}(x;\lambda)
          \;=\; \sum_{\ell \ge 0} \lambda^{\ell}\, x_{t-\ell},
          \label{eq:mmm:adstock}
        \end{equation}
        %
        so a euro spent today still exerts $\lambda^{\ell}$ of its original push $\ell$ weeks
        later. $\lambda = 0.2$ is a one-week echo; $\lambda = 0.8$ is still audible two months out
        (\cref{fig:nb06:curves}, left). The model of \cref{sec:mmm:bayes} truncates the sum at
        $\ell = \nbSixLmax$.

  \item[Saturation, or diminishing returns.] The tenth thousand euros into a channel buys less than
        the first. The response curve used here is the Michaelis--Menten form
        \citep{michaelis1913}, a two-parameter special case of the Hill function \citep{hill1910}
        that has been the workhorse of aggregate advertising response models since
        \citet{little1979}:
        %
        \begin{equation}
          f(a;\ \alpha, \kappa) \;=\; \frac{\alpha\, a}{\kappa + a},
          \label{eq:mmm:mm}
        \end{equation}
        %
        which climbs toward a \emph{ceiling} $\alpha$ --- the most the channel could ever add in a
        week --- and passes half of it at the \emph{half-saturation} spend $a = \kappa$
        (\cref{fig:nb06:curves}, right).
\end{description}

Sales are then the sum of a baseline, the two saturated channels, a direct seasonal effect and
noise:

\begin{equation}
  \text{sales}_t \;=\; 100
    \;+\; \underbrace{f\big(a_t(x^{\text{tv}};\, 0.4);\ 60,\, 50\big)}_{\text{TV: large, real}}
    \;+\; \underbrace{f\big(a_t(x^{\text{bs}};\, 0.2);\ 8,\, 40\big)}_{\text{brand search: small, real}}
    \;+\; \underbrace{35\, s_t}_{\text{the season, directly}}
    \;+\; u_t,
  \qquad u_t \sim \mathcal N(0,\, \nbSixDgpSigma^2).
  \label{eq:mmm:dgp}
\end{equation}

Read \cref{eq:mmm:dgp} against \cref{eq:mmm:xbs} and the confound is complete: the season lifts
sales directly ($35\,s_t$) \emph{and} drives brand-search spend. Brand search is a genuine but
\emph{small} channel riding a large wave it did not make.

Summed over the window, the planted incremental contributions are \eur{\nbSixTrueTv} for
television and \eur{\nbSixTrueBs} for brand search --- a ratio of \nbSixTrueRatio{} to one. In
return-on-spend terms the truth is \nbSixTrueRoasTv$\times$ for television and
\nbSixTrueRoasBs$\times$ for brand search. Note that \emph{both are below one}: on the revenue
definition of ROAS used throughout this chapter, neither channel pays for itself even at a
hundred per cent margin. That fact is the grade every estimator below is marked against, and it is
the last time in the chapter that anything is known for certain.

\input{build/floats/nb06_confound}
\input{build/floats/nb06_curves}

% =====================================================================================
\section{Identification}
\label{sec:mmm:ident}

\subsection*{The estimand}

Let $x_{c,t}$ be spend in channel $c$ in week $t$, and $Y_t\big(x_{c}, x_{-c}\big)$ the potential
sales path under an intervention that sets channel $c$'s spend to a given schedule while leaving
every other channel at its observed one \citep{rubin1974, pearl2009}. The \emph{incremental
contribution} of channel $c$ is what would vanish if the channel were switched off:

\begin{equation}
  \tau_c \;=\; \sum_{t}
    \Big(\ \E\big[\, Y_t \;\big|\; do(x_{c} = x_{c}^{\text{obs}})\,\big]
      \;-\; \E\big[\, Y_t \;\big|\; do(x_{c} = 0)\,\big]\ \Big),
  \label{eq:mmm:estimand}
\end{equation}

with $do(\cdot)$ denoting an intervention --- setting spend by decree rather than merely observing
it. Dividing by what the channel cost gives the quantity the budget meeting actually speaks in:

\begin{equation}
  \text{ROAS}_c \;=\; \frac{\tau_c}{\sum_t x_{c,t}} .
  \label{eq:mmm:roas}
\end{equation}

Two warnings travel with \cref{eq:mmm:roas} and neither is pedantic. First, it is \emph{revenue}
per euro, not profit: break-even at $1\times$ assumes a hundred per cent gross margin, and at a
realistic margin $m$ the break-even bar is $1/m$ --- at $m = \pc{\nbSixMarginThirty}$,
\nbSixBreakevenThirty$\times$. Second, it is an \emph{average}, and budgets are not set at the
average. \Cref{sec:mmm:euro} returns to both.

\subsection*{The causal graph}

The graph of \crefrange{eq:mmm:xtv}{eq:mmm:dgp} is small and completely explicit:
$\text{season} \to \text{sales}$, $\text{season} \to \text{brand search}$,
$\text{TV} \to \text{sales}$, $\text{brand search} \to \text{sales}$. Season is a common cause of
brand-search spend and of sales, so the path
$\text{brand search} \leftarrow \text{season} \to \text{sales}$ is an open backdoor, and Pearl's
backdoor criterion \citep{pearl1995, pearl2009} says: block it by conditioning on season.
Television has no backdoor at all; its arrow into sales is clean.

So the adjustment set is $\{\text{season}\}$, it is known exactly, it is measured exactly, and it
is in every model fitted below. Hold that thought. It is going to fail to be enough, and the
failure is the chapter.

\subsection*{The assumptions}

\begin{assumption}[No unobserved driver correlated with spend]\label{as:mmm:noconf}
Conditional on the controls in the model, the errors $u_t$ of \cref{eq:mmm:dgp} are independent of
every channel's spend: there is no unmeasured variable that moves sales \emph{and} moves the media
plan. Formally, $\{\,x_{c,t}\,\}_{c} \perp u_t \mid \text{controls}$.
\end{assumption}

This is the assumption on which media mix modelling stands, and it deserves to be stated as
harshly as it is true. It is \textbf{untestable} --- an omitted variable is, by definition, not in
the dataset that would test for it --- and in every real marketing dataset it is also
\textbf{false}. Promotions are timed to demand. Competitor activity is unmeasured and correlated
with one's own. Distribution, pricing and product launches move together with campaigns because
the same planning calendar produces all of them. Seasonality can be controlled for only to the
extent it has been correctly modelled, and it never entirely has been. \Citet{lewis2015} put the
point in the sharpest available form: the return on advertising is so small relative to the
volatility of sales that observational data, at realistic sample sizes, cannot distinguish a good
campaign from a bad one --- their calculation implies experiments orders of magnitude larger than
anything an advertiser runs, and \emph{a fortiori} that a weekly regression on chosen spend cannot
do it either. That paper is the intellectual backbone of this chapter. Anyone who reports a media
mix model's ROAS as though it were a measurement should read it first.

This book's rule --- assumptions get tested or are declared untestable --- forces the conclusion
rather than softening it. \Cref{as:mmm:noconf} cannot be tested, and it is not plausible. What can
be done is to \emph{anchor} the model on something that does not need it, and that is what
\cref{sec:mmm:bayes} does with an experiment.

\begin{assumption}[No reverse causality from sales to spend]\label{as:mmm:noreverse}
Spend causes sales; sales do not cause spend. In the simulator this holds by construction: spend is
drawn in \crefrange{eq:mmm:xtv}{eq:mmm:xbs} before sales exist.
\end{assumption}

In the wild, \cref{as:mmm:noreverse} is exactly what brand search violates, and it is worth being
precise about the mechanism because it is not the same thing as \cref{as:mmm:noconf}. Brand-search
spend is triggered by brand-search \emph{queries}, which are made by people who have already
decided to buy. The channel's spend is therefore a \emph{consequence} of impending sales as much
as a cause of them, and the regression cannot see the difference. \Citet{blake2015} measured it by
turning the channel off; no observational model can. Where a channel is demand-\emph{capturing}
rather than demand-\emph{creating}, the sign of the bias is known in advance: upward, and badly.

\begin{assumption}[Correct response form]\label{as:mmm:form}
Each channel's response is geometric adstock \cref{eq:mmm:adstock} composed with Michaelis--Menten
saturation \cref{eq:mmm:mm}, truncated at $\ell = \nbSixLmax$ weeks.
\end{assumption}

\Cref{as:mmm:form} holds here \emph{exactly}, by construction. This is a deliberate and important
choice: the failure demonstrated in this chapter is therefore \textbf{not} a failure of functional
form. It survives the best case. On real data \cref{as:mmm:form} is another wish, and its
violations compound with everything above.

\begin{assumption}[Sufficient independent variation in spend]\label{as:mmm:overlap}
Each channel's spend varies enough, and independently enough of the other channels and of the
controls, that the data can separate their effects.
\end{assumption}

This one \emph{is} testable, and it is the one that fails. Brand-search spend correlates with the
season at \nbSixCorrBsSeason{}. A linear control for the season and a flexibly-saturated
brand-search curve are, on this data, close to the same regressor --- and a likelihood cannot
choose between two things it cannot tell apart. \Cref{sec:mmm:valid} measures precisely how little
it can.

\begin{assumption}[Stable parameters]\label{as:mmm:stable}
The decay $\lambda_c$, ceiling $\alpha_c$ and half-saturation $\kappa_c$ are constant over the
\nbSixNWeeks{} weeks. Creative quality, competitive intensity and audience composition do not
drift.
\end{assumption}

Untestable on one window, and false over any window long enough to fit a media mix model on. This
is the tension every practitioner lives with: \citet{lewis2015}'s arithmetic wants \emph{more}
weeks, and \cref{as:mmm:stable} wants \emph{fewer}. Time-varying-coefficient models
\citep{ng2021} exist precisely to negotiate it.

% =====================================================================================
\section{The classical baseline}
\label{sec:mmm:classical}

Before any sampler: what does a competent analyst without one produce? Most agencies produce
exactly this, and it is worth doing well, because its \emph{specific} failure motivates everything
after it. The estimand is unchanged, and so is the identification argument --- the backdoor set is
$\{\text{season}\}$, so the season enters the regression as a control. The causal work lives in
the graph, not in the estimator. Only the estimation changes.

\subsection*{The obstruction, and the way round it}

Write the media mix model as a regression. The obstruction is immediate: it is not linear in
$(\lambda_c, \kappa_c)$, which sit \emph{inside} the transforms. But it is linear in everything
else, and the algebra says so. Divide \cref{eq:mmm:mm} by its own ceiling and what is left is a
\emph{unit-ceiling basis},

\begin{equation}
  g_{c,t}(\lambda_c, \kappa_c)
    \;=\; \frac{a_{c,t}(\lambda_c)}{\kappa_c + a_{c,t}(\lambda_c)} \;\in\; (0, 1),
  \qquad\text{so that}\qquad
  \text{sales}_t \;=\; \beta_0 \;+\; \sum_c \alpha_c\, g_{c,t} \;+\; \gamma\, s_t \;+\; u_t .
  \label{eq:mmm:ols}
\end{equation}

The ceiling $\alpha_c$ comes back out as the ordinary least-squares \emph{coefficient} on that
basis. So \emph{given} $(\lambda_c, \kappa_c)$ the model is a plain linear regression, and the
classical recipe writes itself: profile. Search a grid of transform values, fit exact least
squares at each, keep the one with the smallest residual sum of squares, and report the regression
at the winner.

That is what is done here, over \nbSixClNGrid{} $(\lambda, \kappa)$ settings --- $\lambda$ across
$[0, 0.8]$ and $\kappa$ from $0.2\times$ to $5\times$ each channel's own mean weekly spend, the
band a practitioner would call plausible. The winner is
$\hat\lambda_{\text{tv}} = \nbSixClLamTv$, $\hat\kappa_{\text{tv}} = \eur{\nbSixClKapTv}$,
$\hat\lambda_{\text{bs}} = \nbSixClLamBs$, $\hat\kappa_{\text{bs}} = \eur{\nbSixClKapBs}$
(the planted truth is $0.40 / 50$ and $0.20 / 40$), at $R^2 = \nbSixClRsq$.

\subsection*{The standard errors, and what ``correct'' means}

Weekly sales are serially correlated, and the residuals of \cref{eq:mmm:ols} carry a lag-one
autocorrelation of \nbSixClAcfOne. The independent-errors variance formula therefore does not
apply, and the intervals are built with Newey--West heteroskedasticity- and
autocorrelation-consistent standard errors \citep{newey1987} at
$\text{maxlags} = \nbSixClMaxlags$ weeks.

One result here corrects a piece of folklore worth correcting. The HAC-to-independent standard
error ratio comes out at $\times\nbSixClSeInflTv$ for television and $\times\nbSixClSeInflBs$ for
brand search --- that is, HAC came out \emph{tighter} for one channel and \emph{wider} for the
other. A HAC covariance \textbf{re-weights}; it does not merely inflate. The point of using it is
correctness, not conservatism, and an analyst who reaches for it expecting a uniformly wider
interval has misunderstood what it does.

\subsection*{The result, and the grade}

\input{build/tables/nb06_compare}

The profiled classical model (\cref{tab:nb06:compare}, fourth row) credits television with
\eur{\nbSixClContribTv} [90\,\%: \eur{\nbSixClContribTvLo}, \eur{\nbSixClContribTvHi}] and brand
search with \eur{\nbSixClContribBs} [\eur{\nbSixClContribBsLo}, \eur{\nbSixClContribBsHi}]. Against
the planted \eur{\nbSixTrueTv} and \eur{\nbSixTrueBs}: the \emph{levels} are off by factors of
\nbSixClRatioTv{} and \nbSixClRatioBs{} respectively, but the \textbf{ranking is
\nbSixClRanking} --- correct. Remember that. \Cref{sec:mmm:compare} has to account for it.

\subsection*{The interval it reports, and the interval it owes}

Now look again at the four numbers above the table. $\hat\lambda$ and $\hat\kappa$ came out of a
grid search and carry \textbf{no interval at all}. The regression that produced those neat
confidence intervals was handed them \emph{as if they were known constants}. So the question a
careful reader should ask is not ``is the point estimate good?'' but: \emph{how much would the
answer move if the transform had been fixed somewhere else --- somewhere the data cannot rule
out?}

The classical fit is structurally unable to answer that, because it conditioned on the answer. So
the question is asked from outside, using the standard Gaussian profile-likelihood result: the set
of $(\lambda, \kappa)$ the data cannot reject at the 90\,\% level is
$\{\, \text{SSE}(\lambda,\kappa) \le \text{SSE}(\hat\lambda,\hat\kappa)
\exp(\chi^2_{4,\,0.90} / n) \,\}$.

That region holds \nbSixRegN{} of the \nbSixClNGrid{} searched settings --- \pc{\nbSixRegShare} of
the entire grid. Inside it, $\lambda_{\text{tv}}$ spans
$[\nbSixRegLamTvLo, \nbSixRegLamTvHi]$ and $\lambda_{\text{bs}}$ spans
$[\nbSixRegLamBsLo, \nbSixRegLamBsHi]$ --- essentially the whole searched band. The decay and the
half-saturation point are, in practice, \textbf{not identified}.

\input{build/floats/nb06_profile}

\Cref{fig:nb06:profile} is what that costs. Every dot is a defensible media mix model. The fit
barely separates them: $R^2$ varies in the fourth decimal across settings whose implied television
ROAS differs several-fold. The implied TV ROAS across the region spans
$[\nbSixEnvLo, \nbSixEnvHi]$ --- width \nbSixEnvWidth{} --- while the interval the model actually
\emph{reports} is $[\nbSixCiLo, \nbSixCiHi]$, width \nbSixCiWidth. The reported interval is
\textbf{\nbSixEnvOverCi$\times$ narrower} than the ignorance it is sitting on. The true TV ROAS of
\nbSixTrueRoasTv$\times$ falls \nbSixClTruthInCi{} the reported interval and
\nbSixClTruthInEnv{} the envelope.

\input{build/tables/nb06_transform}

And the spreadsheet version of the same disease is worse, because it is what most inherited models
actually do. \Cref{tab:nb06:transform} simply \emph{asserts} the transform --- ``we use
industry-standard values'' --- across \nbSixAssertN{} plausible settings. The fits are
indistinguishable: $R^2$ spans only \nbSixAssertRsqLo{} to \nbSixAssertRsqHi. The answers are not:
the implied television ROAS spans \nbSixAssertRoasTvLo{} to \nbSixAssertRoasTvHi, and
\nbSixAssertNInverted{} of the \nbSixAssertN{} settings invert the channel ranking outright.
\textbf{The choice, not the data, is setting the answer.}

\begin{remark}[If you inherit a model whose adstock was ``set to industry standards'']
Its ROAS intervals are decorative. They are correct answers to the question \emph{given that
$\lambda$ and $\kappa$ are exactly these numbers, how precise is the ceiling?} --- a question
nobody asked, conditioned on constants nobody knows. Fixing a parameter does not remove its
uncertainty. It hides it.
\end{remark}

% =====================================================================================
\section{The Bayesian model}
\label{sec:mmm:bayes}

The Bayesian media mix model \citep{jin2017} puts priors on $\lambda_c$, $\alpha_c$ and $\kappa_c$
and \emph{integrates over them} instead of fixing them by fiat. That is exactly the ignorance
\cref{fig:nb06:profile} just measured, so the move is not decoration --- it is the principled
repair of the previous section's central defect. The model, fitted with
\code{pymc-marketing}, is

\begin{align}
  \text{sales}_t \mid \theta
    &\;\sim\; \mathcal N\Big(\beta_0 \;+\; \textstyle\sum_c f_c\big(a_{c,t}\big)
                             \;+\; \gamma\, s_t,\ \ \sigma^2\Big),
      \label{eq:mmm:lik}\\
  a_{c,t} &\;\propto\; \textstyle\sum_{\ell=0}^{\nbSixLmax - 1} \lambda_c^{\ell}\, x_{c,t-\ell},
  \qquad
  f_c(a) \;=\; \frac{\alpha_c\, a}{\kappa_c + a},
      \label{eq:mmm:transforms}
\end{align}

with the same backdoor control $s_t$ as \cref{eq:mmm:ols} --- the graph is handed to the model
explicitly, as a DAG with \code{treatment\_nodes} and an \code{outcome\_node}, so the adjustment
set is derived rather than assumed. Inference is by the No-U-Turn sampler \citep{hoffman2014} in
PyMC \citep{salvatier2016}: \nbSixNDraws{} post-warmup draws across \nbSixNChains{} chains at
$\text{target\_accept} = \nbSixTargetAccept$ --- the saturation geometry needs it.

The model is fitted \textbf{twice}, on identical data, differing in exactly one prior.

\subsection*{(a) Default priors}

Library defaults throughout; in particular the saturation ceilings get
$\alpha_c \sim \text{Gamma}(\mu = 2,\ \sigma = 1)$ in the sampler's max-abs-scaled units, which is
an implied ceiling near \emph{twice} peak weekly sales --- effectively unconstrained. This is the
honest out-of-the-box model: the correct graph, the correct functional form, the correct control,
and no thumb on any scale.

It converges cleanly: $\hat R = \nbSixDefRhat$, minimum bulk effective sample size \nbSixDefEss{}
\citep{vehtari2021}, \nbSixDefDiv{} divergent transitions. And it credits television with
\eur{\nbSixDefTv} and brand search with \eur{\nbSixDefBs}, against a planted \eur{\nbSixTrueTv}
and \eur{\nbSixTrueBs}.

\textbf{It ranks brand search above television.} It over-credits brand search by a factor of
\nbSixDefBsOverTrue. It is exactly, expensively backwards --- and it is the most widely deployed
model in this book.

\subsection*{Why it inverts}

Three mechanisms, and each is checkable rather than asserted.

\emph{One: the collinearity.} Brand-search spend correlates with the season at \nbSixCorrBsSeason.
The \emph{linear} control $\gamma\,s_t$ and the \emph{flexibly saturated} brand-search curve
$f_{\text{bs}}(a_{\text{bs},t})$ are close to the same regressor in this sample
(\cref{as:mmm:overlap} fails).

\emph{Two: the likelihood cannot choose.} Two decompositions --- one crediting the wave to the
season, one crediting it to brand search --- fit the observed series equally well.
\Cref{sec:mmm:valid} measures this, and it is decisive.

\emph{Three: the level account breaks.} Somebody has to pay for the stolen credit, and it is the
baseline. The default fit's intercept lands at \eur{\nbSixDefBaseline} against a planted
\eur{\nbSixTrueBaseline} --- the classic large negative baseline of a non-identified media mix
model. The fitted brand-search ceiling is \eur{\nbSixDefCeilingBs} per week against a planted
\eur{\nbSixTrueCeilingBs}: inflated \nbSixDefCeilingBsOverTrue-fold.

A correct backdoor adjustment was \emph{necessary} and was \emph{not sufficient}. A linear
confounder control does not de-confound a channel whose spend is collinear with the confounder and
which enters through a flexible functional form. That sentence is the chapter, and no prior fixes
it --- because the problem is not that the model believes the wrong thing. The problem is that the
data cannot tell it which thing to believe.

\subsection*{(b) The repair: a prior calibrated from an experiment}

Since the likelihood is silent, the information must come from outside it. The standard and correct
move --- Google's \citep{chan2017} and Uber's \citep{ng2021} recommendation alike, and the one
this book's Chapter~9 exists to supply --- is to \textbf{run an experiment and calibrate the model
against it}. A geo lift test randomizes a channel's spend across matched groups of regions and
reads the sales gap as that channel's true incremental lift, uncontaminated by the seasonality no
regression can net out. That is a \emph{design}, not a model, and it does not need
\cref{as:mmm:noconf}.

The translation from experimental euros to a prior is explicit. The test delivers brand search's
total incremental sales over its window, $\hat L \pm \text{SE}$. The model says brand search's
weekly contribution at its observed operating point is
$f(\bar a) = \alpha_{\text{bs}}\, \bar a / (\kappa_{\text{bs}} + \bar a)$. Matching the two pins
the ceiling:

\begin{equation}
  \alpha_{\text{bs}}
    \;\approx\; \underbrace{\frac{\hat L}{n_{\text{weeks}}}}_{\text{the weekly lift the test measured}}
    \cdot \underbrace{\frac{\kappa_{\text{bs}} + \bar a}{\bar a}}_{\text{ceiling, from the operating point}},
  \qquad
  \alpha^{(s)}_{\text{bs}} \;\sim\; \text{Gamma}\big(\mu = 0.3,\ \sigma = 0.15\big),
  \label{eq:mmm:calib}
\end{equation}

with $\sigma$ scaled from the experiment's own standard error --- wide enough that the data may
still disagree, tight enough to break the ridge. Television, whose incremental role no experiment
has pinned down, keeps its default prior. This is not a thumb on the scale. It is using
experimental evidence at exactly the one place observational data is silent, and
\cref{sec:mmm:valid} sweeps the prior to prove the claim rather than assert it.

The calibrated fit converges as cleanly as the default ($\hat R = \nbSixCalRhat$, minimum bulk ESS
\nbSixCalEss, \nbSixCalDiv{} divergences) and credits television with \eur{\nbSixCalTv} and brand
search with \eur{\nbSixCalBs}. \textbf{The ranking is restored} (\cref{tab:nb06:compare}, last
row). The brand-search ceiling comes back to \eur{\nbSixCalCeilingBs} against the planted
\eur{\nbSixTrueCeilingBs}, and the baseline to \eur{\nbSixCalBaseline} against \eur{100}.

\begin{remark}[In production, feed the experiment to the likelihood, not the prior]
\code{pymc-marketing} can ingest lift tests directly, appending each measurement to the likelihood
as a penalty on the saturation curve's rise over the test's spend delta. That is cleaner than
hand-translating one experiment into one Gamma prior, it accepts many tests at once, and it carries
each test's standard error natively. The single-prior version is used here because it makes the
mechanism visible.
\end{remark}

% =====================================================================================
\section{Diagnostics and validation}
\label{sec:mmm:valid}

\subsection*{The posterior predictive check that cannot save you}

The standard model-criticism question is whether the model can reproduce the data it was trained
on. Both fits are asked it: draw replicated sales series from the posterior predictive, lay the
5--95\,\% band over the observed series, and count the weeks inside.

\input{build/floats/nb06_ppc}

\Cref{fig:nb06:ppc} is the most uncomfortable figure in this chapter. The default fit --- which
credits brand search \eur{\nbSixDefBs} against a planted \eur{\nbSixTrueBs} --- puts
\pc{\nbSixPpcCovDef} of weeks inside its band, at a replicate-mean RMSE of \eur{\nbSixPpcRmseDef}
per week. The calibrated fit puts \pc{\nbSixPpcCovCal} inside, at \eur{\nbSixPpcRmseCal}. The
nominal figure is 90\,\%, and the simulator's own noise has standard deviation
\eur{\nbSixDgpSigma}.

\textbf{Two decompositions an order of magnitude apart on brand search are, to the likelihood,
equally good accounts of the same series.} That is the ridge, in the only form the data will ever
show it: not as a diagonal in a scatter of draws, but as two answers with the same score. It is the
same object \cref{fig:nb06:profile} showed classically --- a near-flat likelihood along the
direction that trades brand search's ceiling against the seasonal coefficient --- and it explains
why no amount of posterior predictive checking will catch this failure.

Keep running posterior predictive checks. A model that \emph{cannot} reproduce the series is
disqualified before anything is read off it. But treat them as \textbf{necessary and never
sufficient} for a causal claim. Fit is not identification.

\subsection*{Parameter recovery: where the damage is}

\input{build/tables/nb06_recovery}

\Cref{tab:nb06:recovery} grades the parameters rather than the contributions, and it closes the
loop on the ``not a thumb on the scale'' claim of \cref{sec:mmm:bayes}. The default fit's damage is
\emph{concentrated}: brand search's ceiling --- the one parameter that can impersonate the seasonal
wave --- inflates \nbSixDefCeilingBsOverTrue-fold, and the baseline is dragged to
\eur{\nbSixDefBaseline} to pay for it. The calibrated prior moves \emph{that one number} back
toward the truth and leaves every other parameter to the data. It is surgical, because the
experiment spoke to exactly the quantity the likelihood is silent about.

Note also what calibration does \emph{not} fix. Television's own ceiling and half-saturation stay
wide and level-shifted in \emph{both} fits. Rankings, not levels --- and the caution holds at the
parameter level too.

\subsection*{The prior sweep: ``didn't you just pick the prior you wanted?''}

The sceptic's question, and it deserves a computed answer. Refit the model across a grid of prior
means for brand search's ceiling --- $\mu \in \{0.15,\, 0.3,\, 0.6,\, 1.0,\, 2.0\}$, from half the
calibrated value up to the library default, keeping the Gamma family and $\sigma = \mu/2$
throughout --- and read off where the ranking survives.

\input{build/tables/nb06_sweep}
\input{build/floats/nb06_sweepfig}

\textbf{The honest answer is a bounded one, and it is smaller than the reader is being invited to
hope.} The ranking is correct ($\Prob(\tau_{\text{tv}} > \tau_{\text{bs}}) > 0.5$) only for
$\mu \le \nbSixSweepHoldTo$, and it is \textbf{already flipped at the very next grid point},
$\mu = \nbSixSweepFlipAt$ --- \nbSixSweepNFlipped{} of the \nbSixSweepNGrid{} grid points invert.
The ranking probability falls from \nbSixSweepPLo{} at the tightest prior to \nbSixSweepPHi{} at
the library default.

So state the robustness as the number it is: \textbf{about twofold in the prior mean, and that is
all.} A real geo test would have to have been wrong by roughly a factor of \nbSixSweepFactor{} on
brand search's ceiling --- not by a few per cent, but not by much more than double either ---
before this budget decision changed direction. It does \emph{not} survive the whole
experiment-plausible range, and it does \emph{not} survive to the library default.

The left panel of \cref{fig:nb06:sweepfig} is the honest dial, and its gentleness is itself the
finding: loosening the prior slides brand search's contribution smoothly away from its truth line
and back toward the flat likelihood's preferred answer, with the data pushing back only weakly
anywhere along the way. The likelihood barely resists, so the prior's strength maps almost
one-for-one onto the conclusion. \textbf{That is what ``the experiment, not the estimator, is doing
the identifying'' looks like when you measure it.}

% =====================================================================================
\section{Point estimate versus posterior}
\label{sec:mmm:compare}

Three estimators are now on the table --- the profiled classical model of
\cref{sec:mmm:classical}, the default-prior Bayesian model and the calibrated one --- all targeting
the estimand of \cref{eq:mmm:estimand} under the identical identification argument.
\Cref{tab:nb06:compare} and \cref{tab:nb06:roas} put them side by side.

\input{build/tables/nb06_roas}

\subsection*{The classical arm won, and this chapter says so}

Read the ranking column of \cref{tab:nb06:compare}. On the committed full series, the classical
profiled model ranks television first --- \emph{correctly} --- and the default Bayesian model does
not. The most widely-deployed Bayesian upgrade of the most widely-deployed marketing model made
the decision-relevant answer \textbf{worse} on the quantity a budget is actually set from.

The book does not flatter Bayes here, and the reason it can afford not to is that the diagnosis is
exact. What did the classical arm buy that verdict with? It \emph{fixed} $(\lambda, \kappa)$ by
fiat --- an infinitely strong constraint, defended by nothing --- and in this sample that
constraint happened to forbid the seasonal-wave theft. Run the same profiler on the notebook's
shorter fast-mode series (104 weeks instead of \nbSixNWeeks{}) and \textbf{it inverts the ranking
outright}. Same estimator, same code, a different draw of the world. And
\cref{fig:nb06:profile}'s right-hand panel already showed that the profile region contains settings
that invert. So the classical win is a coin landing the right way up. \emph{That is luck, not
rigour}, and this chapter refuses to run a victory lap on it.

On \emph{levels}, nobody covers themselves in glory. The classical fit and the calibrated posterior
miss the true contributions in \textbf{opposite directions}, by multiples (\cref{tab:nb06:compare}).
Neither is a euro-grade level estimate.

\subsection*{What the posterior bought}

Three things, and the first is the one \cref{sec:mmm:classical} was owed.

\emph{(i) An interval that includes the transform.} The Bayesian model puts priors on $\lambda$ and
$\kappa$ and integrates over them, so its interval propagates precisely the ignorance the classical
confidence interval conditioned away. \Cref{tab:nb06:roas} lines the three up on television's ROAS.
The classical interval is $[\nbSixCiLo, \nbSixCiHi]$ --- and it is narrow because it \emph{assumed
away} the dominant source of uncertainty. Widen it by hand with the profile envelope,
$[\nbSixEnvLo, \nbSixEnvHi]$, and the classical arm lands in the same neighbourhood as the
posterior's $[\nbSixCalRoasTvLo, \nbSixCalRoasTvHi]$. \textbf{The posterior did not add
uncertainty. It declined to hide it.}

\emph{(ii) A native channel for experimental evidence.} A geo lift test enters the Bayesian model
as a prior on the saturation ceiling --- the exact parameter the test measures
(\cref{eq:mmm:calib}) --- or, in production, straight into the likelihood. The classical analogue
is a hard constraint or a ridge penalty on $\alpha_{\text{bs}}$: a prior in disguise, with no
honest interval and no principled way to weigh the test's own standard error against the data.

\emph{(iii) The number the budget meeting needs.}
$\Prob(\tau_{\text{tv}} > \tau_{\text{bs}}) = \nbSixPTvFirst$ is a posterior probability, and it is
compared against an action bar in \cref{sec:mmm:euro}. The classical arm has no such quantity. Its
nearest analogue --- the share of the transform settings the fit cannot rule out that rank
television first, \pc{\nbSixRegRankOk} --- is a statement about a \emph{grid}, not a probability
about the channels, and it cannot be compared to a bar. This is Chapter~2's
proposition on what a confidence interval licenses, arriving in media-planning dress.

\subsection*{What the posterior did not buy}

It did \textbf{not} rescue a non-identified model, and stating this precisely is the most important
paragraph in the chapter.

Chapter~2 established the book's rule: \emph{the prior does almost nothing; the likelihood is
the whole ballgame} --- a twenty-five-fold prior sweep there moved the answer by
\eur{\nbZeroPriorShift}. \textbf{This chapter is the exception that proves the rule, and it proves
it in the direction the rule predicts.} Here the prior does almost \emph{everything}
(\cref{fig:nb06:sweepfig}) --- and the reason is not that Bayes is stronger here. It is that the
likelihood is \emph{flat}. Where the likelihood has nothing to say, whatever else is in the room
decides, and in fit (a) what was in the room was a library default. A generous
$\text{Gamma}(\mu = 2)$ ceiling prior is not a neutral abstention from an opinion. On this data it
is an active invitation to steal the wave.

So the two diagnostics of this chapter --- the flat SSE surface of \cref{fig:nb06:profile} and the
indistinguishable posterior predictive checks of \cref{fig:nb06:ppc} --- are \textbf{the same
object seen twice}. One near-flat likelihood, seen once through a grid search and once through a
sampler. Being Bayesian about it changed the failure's \emph{shape}, not its \emph{existence}, and
with the library's defaults it made it worse.

Priors help only when they carry \emph{real information}. Here exactly one of them did: the geo
experiment's prior on $\alpha_{\text{bs}}$. And the sweep of \cref{sec:mmm:valid} showed the
conclusion tracking that prior's strength almost one-for-one --- which is the honest confession
that \textbf{the experiment, not the estimator, is doing the identifying}. Bayes is what let the
experiment \emph{in}, coherently, with an interval on the other side. It is not what made the model
identified.

\begin{remark}[The consistency with Chapter~2, stated plainly]
Chapter~2 named three load-bearing conditions of the Bernstein--von Mises theorem, and
predicted that this chapter would break the second: \emph{enough data}. It does, and more
precisely than that --- what fails is not the count of weeks but the \emph{information} in them.
\nbSixNWeeks{} weeks in which brand-search spend and the season move together at
\nbSixCorrBsSeason{} contain almost no information about which of the two moves sales. The
likelihood is flat, the prior therefore does not wash out, and the theorem's conclusion --- that
the two arms agree --- correctly fails to hold. The paradigms disagree here \emph{because the
data are uninformative}, which is exactly what the theorem says should happen.
\end{remark}

% =====================================================================================
\section{The decision, in euros}
\label{sec:mmm:euro}

\subsection*{Average versus marginal, which is the whole point}

Budget chases \emph{incremental} return, and it is set at the \emph{margin}. Two different
euros-per-euro numbers get called ROAS, and only the second can price a budget move:

\begin{equation}
  \text{ROAS}_c \;=\; \frac{\sum_t f_c(a_{c,t})}{\sum_t x_{c,t}},
  \qquad\qquad
  \text{mROAS}_c(\bar x) \;=\; \frac{\partial f_c}{\partial x}\bigg|_{x = \bar x}
    \;=\; \frac{\alpha_c\, \kappa_c}{(\kappa_c + \bar x)^2} .
  \label{eq:mmm:mroas}
\end{equation}

On a saturating curve the marginal is \emph{always} below the average, because the slope keeps
falling. In this data television's calibrated average ROAS is \nbSixAvgRoasTv$\times$ and its
marginal ROAS at today's operating spend is \nbSixMargRoasTv$\times$; brand search's marginal is
\nbSixMargRoasBs$\times$. A channel can therefore look excellent on the report card and be worth
almost nothing on the next euro. The optimal split is reached exactly when the marginals equalise,
$\text{mROAS}_{\text{tv}} = \text{mROAS}_{\text{bs}}$, which is what the sweep below finds
numerically at its peak.

\subsection*{The reallocation, priced two ways}

\input{build/floats/nb06_shift}

Move a \pc{\nbSixShiftPct} slice of the brand-search budget --- \eur{\nbSixShiftEur} --- into
television, and price it. \Cref{fig:nb06:shift} does so \emph{saturation-aware}: it walks both
channels along their fitted response curves, refitting those curves on \emph{every posterior draw}
rather than on the posterior mean, so the planning number inherits a genuine 90\,\% band. The
answer is a median gain of \eur{\nbSixSatDelta} [90\,\%: \eur{\nbSixSatDeltaLo},
\eur{\nbSixSatDeltaHi}], with $\Prob(\Delta > 0) = \nbSixSatPPositive$, and the sweep peaks near a
\pc{\nbSixBestShift} shift.

\begin{remark}[Units. These are window totals, not weekly rates.]
\eur{\nbSixSatDelta} is the gain over the \emph{whole \nbSixNWeeks-week planning window}, because
the shift itself, \eur{\nbSixShiftEur}, is \pc{\nbSixShiftPct} of the \emph{whole-window}
brand-search budget. Dividing by \nbSixNWeeks{} to get a weekly figure would be wrong by that same
factor. Every euro figure in this section is a window total. This is not a pedantic footnote: it
is the error most likely to put a wrong number in front of a CFO.
\end{remark}

Now price the identical shift at each posterior draw's \emph{average} ROAS, ignoring saturation
entirely: \eur{\nbSixAvgDelta} [90\,\%: \eur{\nbSixAvgDeltaLo}, \eur{\nbSixAvgDeltaHi}],
$\Prob(\Delta > 0) = \nbSixAvgPPositive$. It is \textbf{\nbSixAvgOverSat$\times$ larger}, and it is
wrong, for a reason \cref{eq:mmm:mroas} states exactly: television is already close to saturated,
so the next euro into it does not buy what the last one did. \textbf{Quote the saturation-aware
median with its band to the CFO. The average-ROAS figure is a direction check, nothing more} ---
by construction its $\Prob(\Delta > 0)$ is just the ranking probability wearing euro clothing.

\subsection*{The rule, and the verdict}

The decision rule is a stated convention applied to the posterior:

\begin{equation}
  \Prob\big(\tau_{\text{tv}} > \tau_{\text{bs}} \mid D\big)\ \begin{cases}
    > 0.8 & \Rightarrow\ \textsc{reallocate},\\
    \text{otherwise} & \Rightarrow\ \textsc{pilot, and test}.
  \end{cases}
  \label{eq:mmm:rule}
\end{equation}

The posterior gives \nbSixPTvFirst{} --- just short of the \nbSixActionBar{} bar. The verdict is
therefore \textsc{pilot}, not \textsc{reallocate}: move a slice, measure it, do not swing the plan.

And then the caution that the rest of this chapter has earned. The \emph{levels} are not
decision-grade. The calibrated model says television's ROAS is \nbSixCalRoasTv$\times$
[\nbSixCalRoasTvLo, \nbSixCalRoasTvHi], overstating the truth of \nbSixTrueRoasTv$\times$ by
\nbSixCalRoasTvOverTrue-fold. Since ROAS is revenue per euro, break-even at a realistic
\pc{\nbSixMarginThirty} gross margin is not $1.0$ but \nbSixBreakevenThirty$\times$ --- which even
the model's own inflated estimate fails, and which the \emph{true} television ROAS of
\nbSixTrueRoasTv$\times$ fails at \emph{any} margin. So the sentence the analysis actually
supports is: \emph{reallocate toward television, because the ranking is trustworthy; and do not
conclude that television pays, because the levels are not.} ``Is television worth it in absolute
euros?'' is a question for the geo experiment of Chapter~9, not for this model.

\input{build/floats/nb06_correction}

\subsection*{What the experiment was worth, in euros}

The most useful number in this chapter is the price of \emph{not} having run the test. Acting on
the \emph{default} model means believing its inverted ranking, and therefore taking the
\textbf{mirror action}: moving \pc{\nbSixShiftPct} of the \emph{television} budget
(\eur{\nbSixMirrorShiftEur}) into brand search. Price both policies at the true channel ROAS the
simulator knows:

\begin{itemize}[leftmargin=*]
  \item act on the \textbf{calibrated} model: \eur{\nbSixTrueDelta} per planning window, gained;
  \item act on the \textbf{default} model: \eur{\nbSixMirrorDelta} per planning window ---
        \emph{value destroyed}.
\end{itemize}

The swing is \eur{\nbSixExperimentSwing} \textbf{per planning window, recurring every window},
and it is a \emph{floor} on what the geo experiment that calibrated the prior was worth, since it
prices only this one \pc{\nbSixShiftPct} decision and nothing else the calibration de-risks. If a
geo test costs less than a few windows of that, it has already paid for itself --- and this is the
argument to take into the room where the measurement budget is set. \Citet{lewis2015} made the case
in principle; \cref{eq:mmm:calib} and this number make it in euros.

% =====================================================================================
\section{What can go wrong}
\label{sec:mmm:wrong}

\subsection*{Everything above, restated as a checklist}

\Cref{fig:nb06:correction} is the chapter in one panel. Naive least squares, omitting the season,
credits brand search \eur{\nbSixNaiveBs} against a truth of \eur{\nbSixTrueBs}. Adding a linear
seasonality control --- the correct backdoor adjustment --- collapses that credit to
\eur{\nbSixAdjBs}. And then the full media mix model, \emph{with the same correct control still in
it}, inverts the ranking. The confounding correction worked in the linear model and failed in the
flexible one, which is the least intuitive result in this book and the one most worth carrying
away: \textbf{flexibility is not free. A functional form flexible enough to absorb the confounder
will.}

\subsection*{The specific traps}

\begin{itemize}[leftmargin=*, itemsep=.4em]
  \item \textbf{The demand-capture channel.} Any channel whose spend is triggered by demand ---
        brand search above all, but also retargeting, and any always-on channel whose budget is
        pacing to traffic --- violates \cref{as:mmm:noreverse} and will be over-credited by any
        observational model. \Citet{blake2015} is the field measurement; \cref{fig:nb06:correction}
        is the simulation. Treat a high measured ROAS on a demand-capture channel as a red flag,
        not a result.

  \item \textbf{A confident model with clean diagnostics.} The default fit's $\hat R$ is
        \nbSixDefRhat, it has \nbSixDefDiv{} divergences, and its posterior predictive check
        covers at \pc{\nbSixPpcCovDef}. \textbf{Every diagnostic looks fine.} Convergence is a
        statement about the sampler, not about identification (Chapter~2).

  \item \textbf{Fixed transform constants.} \Cref{tab:nb06:transform}: \nbSixAssertNInverted{} of
        \nbSixAssertN{} ``industry-standard'' settings invert the ranking, at a spread in $R^2$ of
        \nbSixAssertRsqLo{} to \nbSixAssertRsqHi. If an inherited model's adstock was set by
        convention, its intervals are decorative.

  \item \textbf{Reading levels off any media mix model.} Even after calibration this one
        overstates television's ROAS \nbSixCalRoasTvOverTrue-fold. Rankings survive; levels do
        not. Anchor levels with an experiment.

  \item \textbf{Average ROAS in a budget meeting.} It over-prices the reallocation by
        \nbSixAvgOverSat$\times$ here (\cref{fig:nb06:shift}). Budgets move at the margin.
\end{itemize}

\subsection*{The limit that no method in this chapter removes}

\Cref{as:mmm:noconf} is untestable and false. Everything in \crefrange{sec:mmm:classical}
{sec:mmm:euro} is therefore conditional on an assumption that will not hold on the reader's data,
and no amount of Bayesian machinery, hierarchical structure or functional flexibility repairs it.
The only repair is a \emph{design}: an experiment, which does not need \cref{as:mmm:noconf} because
a coin flip has already broken the dependence between spend and everything unobserved.

That is why this chapter hands off, deliberately and explicitly, to the next one. The geo lift test
of Chapter~9 is not an appendix to media mix modelling. \textbf{It is what makes media mix
modelling admissible.} An uncalibrated media mix model is not a measurement. A calibrated one --- a
model anchored on an experiment, with the experiment carrying the identification and the model
carrying the interpolation between tests --- is the best instrument a marketing organisation has,
and it is the one worth building.

% =====================================================================================
\section{Takeaways}
\label{sec:mmm:take}

\begin{enumerate}[leftmargin=*, itemsep=.45em]
  \item \textbf{The out-of-the-box media mix model inverted the channel ranking, confidently.}
        Given the correct graph, the correct control and the exact functional form that generated
        the data, it credited brand search \eur{\nbSixDefBs} against a planted \eur{\nbSixTrueBs},
        and television \eur{\nbSixDefTv} against \eur{\nbSixTrueTv} --- ranking a demand-capture
        channel above the channel that creates the demand. Its $\hat R$ was \nbSixDefRhat, its
        divergences \nbSixDefDiv, its posterior predictive coverage \pc{\nbSixPpcCovDef}. Nothing
        in its output warned anyone.

  \item \textbf{A correct backdoor adjustment is necessary and not sufficient.} A linear control
        for the confounder does not de-confound a channel whose spend is collinear with it
        ($\rho = \nbSixCorrBsSeason$) and which enters through a flexible functional form. The
        flexible curve re-absorbs the wave the control was meant to remove, and the baseline goes
        to \eur{\nbSixDefBaseline} to pay for it.

  \item \textbf{Fit cannot arbitrate identification.} Two decompositions an order of magnitude
        apart on brand search score \pc{\nbSixPpcCovDef} against \pc{\nbSixPpcCovCal} coverage and
        RMSE \nbSixPpcRmseDef{} against \nbSixPpcRmseCal{} on the same series
        (\cref{fig:nb06:ppc}). The classical version of the same statement is that
        \pc{\nbSixRegShare} of the transform grid survives the 90\,\% profile region
        (\cref{fig:nb06:profile}). One flat likelihood, seen twice.

  \item \textbf{Fixing a parameter hides its uncertainty; it does not remove it.} The classical
        model's reported ROAS interval is \nbSixEnvOverCi$\times$ narrower than the envelope of
        transforms its own fit cannot rule out, and \nbSixAssertNInverted{} of \nbSixAssertN{}
        ``industry-standard'' settings invert the ranking outright (\cref{tab:nb06:transform}).

  \item \textbf{The classical arm won here --- and it won by luck.} On the full series the profiled
        model ranked television first where the default Bayesian model did not, so Bayes made the
        decision-relevant answer \emph{worse} out of the box, and this book says so. But the
        classical arm bought that verdict by fixing $(\lambda, \kappa)$ by fiat, and on the shorter
        series the same estimator inverts. Do not read this as a classical victory lap.

  \item \textbf{The prior did almost everything --- because the likelihood did almost nothing.}
        This is the mirror image of Chapter~2's finding, not a contradiction of it. Where the
        likelihood is flat, whatever else is in the room decides. A default
        $\text{Gamma}(\mu = 2)$ ceiling prior is not a neutral abstention; on this data it is an
        invitation to steal the wave.

  \item \textbf{The repair is an experiment, and its robustness is about twofold.} A geo lift test
        translated into a prior on brand search's ceiling restores the ranking
        (\eur{\nbSixCalTv} against \eur{\nbSixCalBs}). The ranking holds for
        $\mu \le \nbSixSweepHoldTo$ and flips at the next grid point, $\mu = \nbSixSweepFlipAt$ ---
        \nbSixSweepNFlipped{} of \nbSixSweepNGrid{} grid points invert. Quote that twofold margin.
        Do not say ``it survives any reasonable prior''; it does not.

  \item \textbf{Trust the ranking, never the levels.} Even calibrated, the model overstates
        television's ROAS by \nbSixCalRoasTvOverTrue-fold, and the true ROAS of both channels is
        below break-even at any margin. \textbf{And budgets are set at the margin}: television's
        average ROAS is \nbSixAvgRoasTv$\times$ and its marginal ROAS is \nbSixMargRoasTv$\times$,
        so the average-ROAS pricing of the same reallocation over-states it
        \nbSixAvgOverSat$\times$ (\cref{fig:nb06:shift}). The saturation-aware median with its band
        is the number that goes to the CFO.

  \item \textbf{What the test was worth.} Acting on the default model would have moved
        \eur{\nbSixMirrorShiftEur} the wrong way, destroying \eur{\nbSixMirrorDelta} per window
        against the \eur{\nbSixTrueDelta} the calibrated model gains --- a swing of
        \eur{\nbSixExperimentSwing} per window, recurring, and a floor on the value of the geo
        test. \textbf{The experiment is not a luxury item on top of the media mix model. It is what
        makes it admissible} (Chapter~9).
\end{enumerate}

\notebooknote{%
  This chapter is \code{notebooks/06\_incrementality\_mmm.ipynb}, which runs in the legacy
  environment (\code{pymc<6}, kernel \code{cmp-legacy}) because it is built on
  \code{pymc-marketing}. The simulator is \code{cmp.dgp.mmm\_weekly}
  (\nbSixNWeeks{} weeks, seed \nbSixSeed); the classical arm of \cref{sec:mmm:classical} is a
  profile over \nbSixClNGrid{} transform settings with \code{cmp.classical.ols} under a
  Newey--West covariance; the Bayesian arm is \code{pymc\_marketing.mmm.MMM} with
  \code{GeometricAdstock} and \code{MichaelisMentenSaturation}. The notebook runs in a fast
  teaching mode (\code{CMP\_FAST=1}, 104 weeks, short chains) and a full mode
  (\code{CMP\_FAST=0}, \nbSixNWeeks{} weeks); \textbf{every number in this chapter comes from the
  full run}, emitted by \code{cmp.report} and injected here as a macro. The one exception, flagged
  where it appears in \cref{sec:mmm:compare}, is the observation that the classical profiler
  inverts the ranking on the 104-week fast series --- a property of the other mode, reported
  qualitatively and never used as a result. \Citet{jin2017} and \citet{chan2017} are the standard
  references for the Bayesian model and for the case that it must be calibrated against
  experiments; \citet{lewis2015} is the reason.}
